\documentclass[12pt,a4paper]{article}

\usepackage[utf8]{inputenc}
\usepackage[russian]{babel}
\usepackage{geometry}
\usepackage{setspace}

\geometry{margin=2.5cm}
\setstretch{1.5}

\title{Личная эффективность на примере биографий известных людей}
\author{ }
\date{}

\begin{document}

\maketitle

\section*{Введение}

Личная эффективность — это способность человека достигать поставленных целей с наименьшими затратами времени и ресурсов. История знает множество примеров людей, которые благодаря своей дисциплине, мышлению и привычкам достигли выдающихся результатов. Рассмотрим несколько таких личностей и выделим основные принципы их эффективности.

\section{Стив Джобс}

Стив Джобс — сооснователь компании Apple, человек, изменивший индустрию технологий. Его биография показывает, что ключевым фактором успеха является умение фокусироваться.

Джобс уделял внимание только самым важным проектам, отказываясь от второстепенного. Он говорил, что важно не только то, что мы делаем, но и то, от чего мы отказываемся.

\textbf{Черты эффективности:}
\begin{itemize}
\item Фокус на приоритетах
\item Смелость в принятии решений
\item Стремление к качеству
\end{itemize}

\section{Илон Маск}

Илон Маск — предприниматель, стоящий за компаниями Tesla и SpaceX. Его отличает невероятная работоспособность и способность управлять несколькими сложными проектами одновременно.

Он использует метод разбиения задач на маленькие шаги и строго планирует своё время.

\textbf{Черты эффективности:}
\begin{itemize}
\item Тайм-менеджмент
\item Системное мышление
\item Готовность к риску
\end{itemize}

\section{Мария Кюри}

Мария Кюри — выдающийся учёный, лауреат двух Нобелевских премий. Её успех стал возможен благодаря настойчивости и преданности делу.

Несмотря на трудности, она продолжала исследования, не теряя интереса к науке.

\textbf{Черты эффективности:}
\begin{itemize}
\item Упорство
\item Самодисциплина
\item Любовь к своему делу
\end{itemize}

\section*{Правила личной эффективности}

На основе рассмотренных биографий можно сформулировать следующие правила:

\begin{enumerate}
\item Определяйте приоритеты и фокусируйтесь на главном.
\item Разбивайте большие задачи на более мелкие.
\item Управляйте своим временем осознанно.
\item Развивайте дисциплину и настойчивость.
\item Не бойтесь ошибок и воспринимайте их как опыт.
\item Работайте над тем, что вам действительно интересно.
\item Минимизируйте отвлекающие факторы.
\end{enumerate}

\section*{Заключение}

Личная эффективность — это не врождённое качество, а навык, который можно развивать. Примеры известных людей показывают, что успех достигается благодаря сочетанию правильных привычек, мышления и системного подхода к работе.

\end{document}